\section{Zaključak}

Sembot predstavlja inovativno rješenje za automatizaciju podrške studentima na sveučilištima, koristeći moderne tehnologije poput Semantic Routera i FastAPI-ja. Projekt nudi jednostavan, skalabilan i proširiv sustav koji može značajno smanjiti opterećenje na administrativno osoblje, omogućujući studentima brži pristup informacijama. Upotreba semantičkog usmjeravanja omogućuje razumijevanje prirodnog jezika, dok JSON baza podataka pruža fleksibilnost u upravljanju sadržajem.

Trenutna implementacija ima nekoliko prednosti, uključujući jednostavnost postavljanja, intuitivno sučelje i mogućnost proširenja. Međutim, postoje i ograničenja, poput skalabilnosti JSON baze i nedostatka naprednih funkcionalnosti poput sustava preporuke. Buduće nadogradnje uključuju:

\begin{itemize}
    \item Implementaciju sustava preporuke sljedećeg pitanja.
    \item Prijelaz na relacijsku bazu podataka za veću skalabilnost.
    \item Dodavanje višejezične podrške za međunarodne studente.
    \item Razvoj administrativnog sučelja za upravljanje podacima.
\end{itemize}

Integracija naprednijih modela prirodnog jezika, poput lokalno hostiranih modela za smanjenje ovisnosti o vanjskim API-jima.
Sembot ima potencijal postati ključni alat za podršku studentima, s mogućnošću prilagodbe za različite fakultete i obrazovne institucije. Projekt demonstrira snagu modernih tehnologija u rješavanju stvarnih problema i otvara put za daljnje inovacije u području automatizirane korisničke podrške.
